\section{Proposed comparison within developmental stages}
The pooled sample mixes developmental stages. Its dimension--error curve
therefore combines variation within stages with differences between them.
A natural next question is whether the curve is lower or has a different
shape within each stage. This section specifies a proposed experiment;
no stage-specific embeddings have yet been fitted.

There are ten annotated stages. Labels beginning with E denote embryonic
stages; labels beginning with P denote postnatal stages. Table~\ref{tab:stages}
shows the available cells. All stages contain at least 1,500 cells, so a
balanced comparison can draw a new uniform sample of 1,500 cells without
replacement \emph{within each stage}. Splitting the existing pooled
1,500-cell sample would instead give highly unequal, sometimes very small
groups (only 19 E12 cells), which is unsuitable for this comparison.

\begin{table}[htbp]
\centering\begin{tabular}{lrr}\toprule
Stage & Available cells & Proposed sample \\ \midrule
E11 & 9,364 & 1,500 \\
E12 & 1,610 & 1,500 \\
E14 & 26,944 & 1,500 \\
E16 & 5,401 & 1,500 \\
E18 & 23,011 & 1,500 \\
P0 & 9,925 & 1,500 \\
P2 & 17,963 & 1,500 \\
P5 & 6,756 & 1,500 \\
P8 & 11,727 & 1,500 \\
P14 & 8,103 & 1,500 \\
\bottomrule\end{tabular}

\caption{Available developmental-stage annotations in the 120,804-cell
input and the proposed sample size per stage. Counts are cells, not
biological replicates. Stages are ordered developmentally. No stage is
excluded under the proposed balanced sampling rule.}
\label{tab:stages}
\end{table}

For each sampled stage, keep the same full-population 1,814-PC coordinate
system and unscaled Euclidean target metric. Do not refit the input PCs,
whiten coordinates, or standardize each stage separately. The global
99.002235-percent variance retention does not imply that this exact fraction
holds within every stage. A stage-specific neighbor-search PCA retaining
90 percent of that stage sample's variance may still select local
interactions; as before, it must not change target distances.

Fit dimensions $d\in\{3,4,5,10,20,25,50\}$, with one shared preparation
per stage: 200 internally selected distance-proportional landmarks,
ten-neighbor local interactions, regional endpoint counts, and 100 SGD
epochs on a 100-epoch schedule with $\epsilon=0.1$. Use one explicitly
seeded start per stage and dimension for the first pass, with nested
starting-coordinate columns within a stage. Record sampling, preparation
and fitting seeds separately. This gives 70 planned fits, with no
resumption or modification of the paused full-population run.

Evaluate every within-stage unordered pair: 1,124,250 pairs per sample.
Report the all-pairs score and sparse proxy separately, retaining
ten-epoch histories and terminal coordinates without post-fit rescaling.
Show stage-specific curves on common axes, together with the existing
pooled-sample curve over the same dimensions. Retain directly fitted and
remaining-pair diagnostics when the scores disagree. A replenished queue
can use up to ten single-threaded workers under the same 50 GiB aggregate
resident-memory watchdog; preparation and fitting must share both this
memory allowance and the ten-thread CPU budget.

Interpret differences descriptively. The relative score divides each
residual by that stage sample's own target distance; target-distance
distributions can differ across stages. Lower error within a stage would
show easier reconstruction of that sample's Euclidean distances under
this method, not by itself lower biological intrinsic dimension or a
causal effect of age. Cell-type composition and sample or batch structure
may also differ. Equal cell counts control the number of pairs and the
landmark fraction, but do not remove these differences. Multiple starts
and independent cell subsamples should follow before making robust
comparisons between stages.
\FloatBarrier
